% !TEX root = task-2-4-2.tex
% Editable task bank. Append new BankTask blocks at the end of the document.
\documentclass[12pt,a4paper]{article}
\usepackage{physics-textbook}
\usepackage{graphicx,needspace}
\usetikzlibrary{decorations.pathmorphing}
\newif\ifShowSolutions
\ifdefined\TutorEdition
  \ShowSolutionstrue
\else
  \ShowSolutionsfalse % Change to \ShowSolutionstrue to display solutions.
\fi
\newcounter{banktask}
\newsavebox{\BankTaskBox}
\newcommand{\BankTaskContent}[3]{%
  \noindent\textbf{\thebanktask.}\quad #2\par
  \ifShowSolutions
    \Needspace{50mm}% Keep the solution heading with its opening text or diagram.
    #3\par
    \ifdefined\TutorEdition
      \ifcsname TutorSource#1\endcsname
        \nopagebreak\smallskip{\small\csname TutorSource#1\endcsname\par}%
      \fi
    \fi
  \fi
}
\newcommand{\BankTask}[3]{%
  \refstepcounter{banktask}\par\medskip
  \sbox{\BankTaskBox}{\begin{minipage}{\linewidth}%
    \BankTaskContent{#1}{#2}{#3}\end{minipage}}%
  \ifdim\dimexpr\ht\BankTaskBox+\dp\BankTaskBox\relax>\dimexpr\textheight-20mm\relax
    \Needspace{35mm}\BankTaskContent{#1}{#2}{#3}%
  \else
    \Needspace{\dimexpr\ht\BankTaskBox+\dp\BankTaskBox\relax}%
    \noindent\usebox{\BankTaskBox}\par
  \fi
  \medskip
}
\begin{document}
\BookSubsection{2.4.2}{Тепловое равновесие}


% Source sdamgia-3360: https://phys-ege.sdamgia.ru/problem?id=3360
\BankTask{sdamgia-3360}{
Воздух в комнате состоит из смеси газов: водорода, кислорода, азота, водяных паров, углекислого газа и др. При тепловом равновесии у этих газов обязательно одинаковы

$1$) температуры

$2$) парциальные давления

$3$) концентрации молекул

$4$) плотности
}{
\smallskip\noindent\textbf{Решение}\par
При тепловом равновесии газов выравниваются их средние кинетические энергии теплового движения, а значит, и температуры.

Ответ: $1$.
}

% Source sdamgia-6116: https://phys-ege.sdamgia.ru/problem?id=6116
\BankTask{sdamgia-6116}{
В состоянии теплового равновесия в термодинамической системе не прекращается

А) изменение температуры частей этой системы.

Б) обмен энергией между молекулами.

$1$) только А

$2$) только Б 

$3$) и А, и Б

$4$) ни А, ни Б
}{
\smallskip\noindent\textbf{Решение}\par
В состоянии теплового равновесия температура постоянна. Обмен энергией между молекулами всегда происходит внутри вещества за счет их теплового движения.

Правильный ответ указан под номером $2$.
}

% Source sdamgia-3725: https://phys-ege.sdamgia.ru/problem?id=3725
\BankTask{sdamgia-3725}{
В калориметр с водой, имеющей температуру $20^\circ$C, кладут металлический брусок, имеющий температуру $40^\circ$C. Через некоторое время в калориметре устанавливается тепловое равновесие. Как в результате изменятся следующие физические величины: внутренняя энергия бруска, внутренняя энергия воды, суммарная внутренняя энергия системы?

$1$. Увеличится.

$2$. Уменьшится.

$3$. Не изменится.

\begin{center}\begin{tabular}{|p{0.287\linewidth}|p{0.287\linewidth}|p{0.287\linewidth}|}\hline
Внутренняя энергия \newline бруска & Внутренняя энергия \newline воды & Суммарная \newline внутренняя энергия \newline системы\\\hline
\strut & \strut & \strut\\\hline\end{tabular}\end{center}
}{
\smallskip\noindent\textbf{Решение}\par
При установлении теплового равновесия температуры всех тел сравниваются. Следовательно, вода в калориметре нагревается, а брусок остывает. Чем горячее тело, тем больше его внутренняя энергия. Таким образом, внутренняя энергия бруска уменьшается, а внутренняя энергия воды увеличивается. 

Поскольку калориметр изолирует систему, тепло из вне в нее не поступает. Кроме того, никто не совершает работы над ней, а значит, согласно первому началу термодинамики, суммарная внутренняя энергия системы остается неизменной.

Ответ: $213$.
}

% Source sdamgia-36755: https://phys-ege.sdamgia.ru/problem?id=36755
\BankTask{sdamgia-36755}{
\begin{center}\begin{tikzpicture}[task graph]
\fill[blue!8](0,0)rectangle(2.4,1.6);\fill[blue!15](2.4,0)rectangle(4.8,1.6);\draw[task curve](0,0)rectangle(4.8,1.6);\draw[task curve](2.4,0)--(2.4,1.6);\node at(1.2,.8){А};\node at(3.6,.8){Б};
\end{tikzpicture}\end{center}
При изучении явления теплообмена герметичный теплоизолированный сосуд с идеальным одноатомным газом разделили на две одинаковые части неподвижной перегородкой, способной проводить тепло (см. рисунок). В разных частях сосуда газ был нагрет до разных температур. Температура газа в части $A$ была равна $293\,\text{К}$, а в части $\text{Б}$ +$40^\circ$C. Количество газа одинаково в обеих частях. Теплоёмкостью сосуда пренебречь.

$1$. Температура газа в части $\text{Б}$ повысится.

$2$. Внутренняя энергия газа в части $A$ не изменится.

$3$. При теплообмене газ из части $\text{Б}$ отдавал количество теплоты, а газ в части $A$ его получал.

$4$. Через достаточно большой промежуток времени температуры газов в обоих частях станут одинаковыми и равными +$30^\circ$C.

$5$. В результате теплообмена газ, находящийся в части $\text{Б}$, совершил работу.
}{
\smallskip\noindent\textbf{Решение}\par
$1$. Неверно. В начальный момент времени температура газа в части $\text{Б}$ равнялась $313\,\text{К}$. В результате теплообмена температура газа понизится.

$2$. Неверно. При теплообмене температура газа в сосуде $A$ повысится, следовательно, внутренняя энергия изменится.

$3$. Верно. При теплообмене газ из части $\text{Б}$ отдавал количество теплоты, а газ в части $A$ его получал, так как направление теплообмена происходило от $\text{Б}$ к $A$.

$4$. Верно. В результате теплообмена наступит тепловое равновесие, при котором газ в обоих сосудах будет иметь одинаковую температуру. По закону сохранения энергии $U_1 + U_2=U'_1 + U'_2.$ Тогда:

$\frac{3}{2} \nu RT_1 + \frac{3}{2} \nu RT_2= \frac{3}{2} \nu RT + \frac{3}{2} \nu RT,$

откуда $T= \frac{T_1 + T_2}{2} = \frac{293 + 313}{2} =303\text{К}=30 ^\circ\text{C} .$

$5$. Неверно. Поскольку газы не расширялись, то работа ими не совершалась.

Ответ: $34$.
}

% Source sdamgia-12859: https://phys-ege.sdamgia.ru/problem?id=12859
\BankTask{sdamgia-12859}{
В калориметр налит $1\,\text{л}$ воды при температуре $0^\circ$C. В этот калориметр последовательно выливают $50$ одинаковых мензурок воды, нагретой до температуры +$50^\circ$C. Объем мензурки $20\,\text{см}^{3}$. Потерями теплоты и теплоемкостью калориметра можно пренебречь. Какая температура (в °С) установится в калориметре?
}{
\smallskip\noindent\textbf{Решение}\par
Объем добавленной воды $50\cdot 20\,\text{см}^{3}$ = $1000\,\text{см}^{3}$ = $1\,\text{л}$ совпадает с объемом исходной воды, значит, температура установится средней между их температурами, то есть $25^\circ$C.

Ответ: $25$.
}

% Source sdamgia-19663: https://phys-ege.sdamgia.ru/problem?id=19663
\BankTask{sdamgia-19663}{
В калориметр залили три порции воды массами $200\,\text{г}$, $300\,\text{г}$ и $500\,\text{г}$, которые имели температуры $20^\circ$C, $40^\circ$C и $60^\circ$C соответственно. Теплообмен воды с окружающими телами пренебрежимо мал. Какой будет температура воды в калориметре после установления теплового равновесия? (В градусах Цельсия.)
}{
\smallskip\noindent\textbf{Решение}\par
Уравнение теплообмена: $Q_{\text{получ}}=Q_1 + Q_2 + Q_3=0.$ Теплоту вычислим по формуле $Q = cm_x ( t - t_x ) ,$ откуда $Q_{\text{получ}} = cm_1 ( t - t_1 ) + cm_2 ( t - t_2 ) + cm_3 ( t - t_3 ) = 0.$

Раскрываем скобки, приводим подобные и получаем выражение для $t$:

$t = \frac{m_1t_1 + m_2t_2 + m_3t_3}{m_1 + m_2 + m_3} = \frac{0{,}2 \cdot 20 + 0{,}3 \cdot 40 + 0{,}5 \cdot 60}{0{,}2 + 0{,}3 + 0{,}5} = 46 ^\circ\text{C}.$

Ответ: $46$.
}

% Source sdamgia-10943: https://phys-ege.sdamgia.ru/problem?id=10943
\BankTask{sdamgia-10943}{
В калориметр, в котором находилась вода массой $2\,\text{кг}$ при температуре $0^\circ$C, бросили $300\,\text{г}$ льда при температуре -$55^\circ$C. Какая масса льда в граммах окажется в калориметре после установления теплового равновесия? (Удельная теплоемкость льда — $2100\,\text{Дж}$/(кг · °С), удельная теплота плавления льда — $330\,\text{кДж}$/кг.)
}{
\smallskip\noindent\textbf{Решение}\par
Найдем количество теплоты, которое необходимо затратить, чтобы нагреть лед до $0^\circ$C:

$Q=cm\Delta T=2100 \cdot 0{,}3 \cdot 55=34650\text{Дж}.$

Это количество теплоты будет образовано за счет замерзания воды. При этом часть воды превратится в лед массой:

$m= \frac{Q}{\lambda} = \frac{34650}{330000} =105\text{г}.$

Таким образом, после установления теплового равновесия в калориметре окажется $405\,\text{г}$ льда ($300\,\text{г}$ первоначально плюс $105\,\text{г}$ замерзшей воды).

Ответ: $405$.
}

% Source sdamgia-4217: https://phys-ege.sdamgia.ru/problem?id=4217
\BankTask{sdamgia-4217}{
В теплоизолированном цилиндре, разделенном на две части тонким невесомым теплопроводящим поршнем, находится идеальный одноатомный газ. В начальный момент времени поршень закреплен, а параметры состояния газа — давление, объем и температура — в одной части цилиндра равны $p_{1}=1$ атм, $V_{1}=1\,\text{л}$ и $T_{1}=300\,\text{К}$, а в другой, соответственно, $\text{р}_{2}=2$ атм, $V_{2}=1\,\text{л}$ и $T_{2}=600\,\text{К}$. Поршень отпускают, и он начинает двигаться без трения. Какое давление газа установится в цилиндре спустя достаточно долгое время, когда будет достигнуто состояние равновесия? Теплоемкостями цилиндра и поршня можно пренебречь.
}{
\smallskip\noindent\textbf{Решение}\par
Запишем уравнение состояния (уравнение Клапейрона — Менделеева) для газа в обеих частях цилиндра в начальный момент времени:

$p_1V_1=\nu _1RT_1,$ $p_2V_2=\nu _2RT_2.$

Из первого начала термодинамики следует, что внутренняя энергия газа в этом процессе сохраняется, так как газ не обменивается теплотой с окружающими телами и не совершает работы. Запишем выражения для внутренней энергии $U$ газа.

В начальный момент:

$U_1= \frac{3}{2} \nu _1RT_1= \frac{3}{2} p_1V_1,$ $U_2= \frac{3}{2} \nu _2RT_2= \frac{3}{2} p_2V_2,$

а в установившемся состоянии равновесия:

$U= \frac{3}{2} p ( V_1 + V_2 ) =U_1 + U_2= \frac{3}{2} ( p_1V_1 + p_2V_2 ) .$

Отсюда окончательно получаем:

$p=\frac{p_1V_1+p_2V_2}{V_1+V_2}=1{,}5$ атм.

Ответ: $1{,}5$ атм.
}

% Source sdamgia-38954: https://phys-ege.sdamgia.ru/problem?id=38954
\BankTask{sdamgia-38954}{
Теплоизолированный цилиндр, расположенный горизонтально, разделён подвижным теплопроводящим поршнем на две части. В одной части цилиндра находится гелий, а в другой аргон. В начальный момент температура гелия равна $600\,\text{К}$, а аргона — $900\,\text{К}$, объемы, занимаемые газами, одинаковы, а поршень находится в равновесии. Во сколько раз изменится объём, занимаемый гелием, после установления теплового равновесия, если поршень перемещается без трения? Теплоёмкостью цилиндра и поршня пренебречь.
}{
\smallskip\noindent\textbf{Решение}\par
По условию первоначально поршень находился в равновесии. Следовательно, давления гелия и аргона были одинаковыми, которые обозначим $p_{0}$. При этом газы занимали одинаковый объем $V$. Запишем уравнение Клапейрона-Менделеева для каждого газа:

$p_0V=\nu_1RT_1,$

$p_0V=\nu_2RT_2,$

где $\nu_1$ и $\nu_2$ — количества вещества гелия и аргона. Тогда получаем соотношение:

$\nu_1T_1=\nu_2T_2 ( 1 ) .$ 

Поршень является теплопроводящим, поэтому через некоторое время в обоих сосудах установится одинаковая температура $T$. При отсутствии силы трения поршень придет в движение и остановится, когда давления газов станут снова одинаковыми и примут некоторое значение $p$, а объемы газов станут равными $V_{1}$ и $V_{2}$ соответственно. Запишем уравнение Клапейрона-Менделеева для нового состояния газа:

$pV_1=\nu_1RT,$

$pV_2=\nu_2RT.$

Тогда получим соотношение:

$\frac{V_2}{V_1} = \frac{\nu_2}{\nu_1} ( 2 ) .$

Объединяя уравнения ($1$) и ($2$), а также учитывая, что $V_1 + V_2=2V,$ получим:

$\frac{V_1}{V} = \frac{2T_2}{T_1 + T_2} = \frac{2 \cdot 900}{600 + 900} =1{,}2.$

Таким образом, объем, занимаемый гелием, увеличится в $1{,}2$ раза.

Ответ: увеличится в $1{,}2$ раза.
}

% Source sdamgia-9757: https://phys-ege.sdamgia.ru/problem?id=9757
\BankTask{sdamgia-9757}{
Для получения и поддержания температуры $0^\circ$С, одной из двух реперных точек на шкале Цельсия, в лабораторной практике часто используют следующий метод. В теплоизолированный стакан наливают дистиллированную воду комнатной температуры, поливают воду сверху жидким азотом, перемешивая смесь ложкой до тех пор, пока не образуется масса серого цвета, состоящая из мелких кристалликов льда и воды. Это обеспечивает нужную температуру в течение длительного времени — смесь помещают в сосуд Дьюара, где она медленно тает при $0^\circ$С. Какой объем $V$ жидкого азота требуется израсходовать для получения массы $m=300\,\text{г}$ такой смеси, содержащей $75$\% льда и $25$\% воды (по массе), из воды при $25^\circ$С? Теплоемкостями стакана и ложки, а также потерями теплоты можно пренебречь. Плотность жидкого азота $\rho _{\text{ж}}=808\,\text{кг}$/$\text{м}^{3}$, удельная теплота парообразования $L=197{,}6\,\text{кДж}$/кг.
}{
\smallskip\noindent\textbf{Решение}\par
$1$. Из условия задачи понятно, что происходит в системе при получении смеси льда и воды, реализующей реперную точку шкалы Цельсия — $0^\circ$С. Сначала вся масса воды с удельной теплоемкостью $c_{\text{в}}$ охлаждается от $25^\circ$С до $0^\circ$С на $\Delta T=25^\circ$С, а затем $75$\% массы воды превращается в лед при температуре $0^\circ$С, отдавая свою теплоту кристаллизации (удельная теплота кристаллизации воды $\lambda$ равна удельной теплоте плавления льда). В этом процессе отданное водой количество теплоты равно:

$Q = c_{\text{в}} \cdot m \cdot \Delta T + \lambda \cdot ( 3m/4 ) = 4200 \cdot 0{,}3 \cdot 25 + 3{,}3 \cdot 10 ^{5} \cdot 0{,}225 = 105{,}75\,\text{кДж}$.

$2$. Теплота $Q$, отбираемая от воды в данном процессе, идет по условию только на испарение массы $m_{\text{ж}}=\rho _{\text{ж}}\cdot V$ жидкого азота: $Q=mL=\rho_{\text{ж}}VL=808 \cdot 197600V=159660{,}8V\text{кДж}.$

$3$. Таким образом, надо израсходовать объем жидкого азота, равный 

$V=105{,}75$/$159660{,}8$ $\approx$ $6{,}623\cdot 10^{-4}$ $\text{м}^{3}$ $\approx$ $662\,\text{см}^{3}$.

Ответ: $V$ $\approx$ $662\,\text{см}^{3}$.
}


% Source 0443DE: https://ege.fipi.ru/bank/questions.php?search=1&proj=BA1F39653304A5B041B656915DC36B38&qid=0443DE
\BankTask{0443DE}{
Герметичный теплоизолированный сосуд разделили неподвижной перегородкой, способной проводить тепло, на две равные части. В первую часть сосуда поместили некоторое количества аргона при температуре \(328\,\text{К}\), а во вторую --- такое же количество аргона при температуре \(15\,^{\circ}\mathrm{C}\).
Теплоёмкость сосуда пренебрежимо мала.
Рассматриваются изменения, происходящие с аргоном при переходе к тепловому равновесию.
\begin{enumerate}
\item Температура газа во второй части сосуда повысилась.
\item Внутренняя энергия газа в первой части сосуда увеличилась.
\item При теплообмене газ в первой части сосуда отдавал теплоту, а газ во второй части сосуда её получал.
\item Через достаточно большой промежуток времени температура газов в обеих частях сосуда стала одинаковой и равной \(25\,^{\circ}\mathrm{C}\).
\item В результате теплообмена газ в первой части сосуда совершил положительную работу.
\end{enumerate}
}{
\smallskip\noindent\textbf{Решение от GPT}\par
Начальная температура во второй части \(T_2=15+273=288\,\text{К}\). Объёмы постоянны, поэтому работа каждого газа \(A=0\). По закону сохранения энергии для равных количеств одного газа
\[
T_{\text{равн}}=\frac{328+288}{2}=308\,\text{К}=35\,^{\circ}\mathrm{C}.
\]
\begin{enumerate}
\item Верно: температура второго газа повысилась от \(288\) до \(308\,\text{К}\).
\item Неверно: \(U=\frac32\nu RT\), поэтому при охлаждении первого газа его внутренняя энергия уменьшилась.
\item Верно: теплота передавалась от более горячего первого газа ко второму.
\item Неверно: равновесная температура равна \(35\,^{\circ}\mathrm{C}\).
\item Неверно: при неподвижной перегородке \(A=0\).
\end{enumerate}
Ответ: \(13\).
}

\end{document}
